\scanpage{52}
\deleted{\unsure{dont}} de degré $3$, dont les \inserted{$3$} fibres sont les \inserted{trois} triades (\unsure{mutuellmt|mutuellt|mutuellement} disjointes) formant l'une des deux familles de triades de $X_\beta$. Alors

a) Les \deleted{\unsure{autres}} deux familles de triades, fibres des deux applications $X_\alpha\to I_\beta$, $X_\alpha\unsure{-|\to}I_\gamma$ (où \deleted{\unsure{$\alpha$}} $A=\{\alpha,\beta,\gamma\}$) sont distinctes, donc l'application correspondante $X_\alpha\xrightarrow{\sim}I_\beta\times I_\gamma$ est bijective.

Plus intrinsèquemt, $\forall\alpha\in A$, on a
\[
 X_\alpha\xrightarrow{\sim}\prod_{\beta\in I-\{\alpha\}} I_\beta
\]
d'où une bijection naturelle
\[
 S\simeq\coprod_\alpha\prod_{\beta\in I-\{\alpha\}}I_\beta\quad\deleted{\unsure{I_{\alpha}}}
\]
dans le cas où $A=\{1,2,3\}$, cela fait
\[
 S\simeq\underbrace{I_2\times I_3}_{X_1}\ \amalg\
 \underbrace{I_3\times I_1}_{X_2}\ \amalg\
 \underbrace{I_1\times I_2}_{X_3}
\]
On a $\mathrm{coor}(X_\alpha)\simeq\deleted{\unsure{I_\alpha}}A-\{\alpha\}$.

c) Pour que \deleted{$x_\alpha\in X_\alpha,y_\alpha\in$ \unsure{avec} $\unsure{x\ne y}$} $x_\alpha,y_\alpha\in X_\alpha$ \deleted{$\unsure{x_\alpha\ne y_\alpha}$}, soient \uline{non liés} il f. et s. que $\exists\beta\in A-\{\unsure{\alpha|\beta}\}$ \deleted{\unsure{que}} tel que $p_{\beta\alpha}(x_\alpha)=p_\beta(y_\alpha)$ (si $x_\alpha\ne y_\alpha$, ce $\beta$ est unique) -- pour qu'ils soient liés, il f et s. que $\forall\beta\in A-\{\alpha\}$, on ait $p_{\beta\alpha}(x_\alpha)\ne p_{\beta\alpha}(y_\alpha)$

d) Pour que $x_\alpha\in X_\alpha$ et $y_\beta\in X_\beta$ $(\alpha\ne\beta)$ soient liés il f et s que (si $\gamma$ \deleted{tel que} est l'élémt $\ne\alpha,\beta$ de $A$) on ait $p_{\gamma\alpha}(x_\alpha)=p_{\gamma\beta}(x_\beta)$
% archive sheet number: 51
